\section{Discussion}

\subsection{Why the inversion exists}
\label{sec:why}

The inversion in Fig.~\ref{fig:inversion} invites a tempting misreading: that measurement
reveals bench work to be safer than feared, and that the unmeasured steps are
over-scored by an instrument reflecting folklore. I argue the opposite, on three grounds.

First, the measurement that exists was not allocated by hazard. It was allocated by
occasion. The pre-analytical measurements were performed during the COVID-19 pandemic, in
response to acute concern among biochemistry staff about handling specimens from infected
patients~\cite{dubuis2021}. The cartridge-assay measurement was performed to support
deployment of a new diagnostic platform into settings lacking containment
infrastructure~\cite{banada2010}. Neither study was commissioned because the step in
question was judged the most dangerous in the laboratory. They were commissioned because
somebody needed an answer about that step at that moment. Microbiology bench work attracted
no comparable occasion, and so attracted no comparable measurement.

Second, the unmeasured steps are not unevidenced. They are evidenced by a different and
less tractable method: counting the people who became infected. An occupational attack
rate roughly sixty-five times background, with a case fatality of eight in
sixteen~\cite{sejvar2005}, is a strong signal about a procedure. It is simply not a signal
expressed in CFU/m\textsuperscript{3}, and an evidence framework that recognises only
aerosol measurement would discard it. This is the reason the grading scale treats incident
evidence as its own level rather than as missing data.

Third, the direction of the gap is itself informative about where measurement effort has
value. The steps we have measured turn out to release little; the steps we have not
measured are those associated with documented infection. If the purpose of measurement is
to reduce uncertainty where the stakes are highest, the existing allocation is close to
inverted, and correcting it is inexpensive relative to the question.

\subsection{What a negative result licenses}

The negative pre-analytical findings require careful statement, and an earlier draft of
this work stated them too strongly.

Both negative results were obtained \emph{with containment in place}. Centrifugation was
performed with closed-cap tubes in a swing-bucket rotor; tube opening was performed behind
a plexiglass panel~\cite{dubuis2021}. The finding is therefore evidence that those
controls work. It is not evidence that the bare open bench is adequate, and a laboratory
citing this study to justify removing a barrier would be inverting its meaning.

The detection limit compounds this. At 4~CFU/m\textsuperscript{3}, release below that
threshold would be reported as zero. For an organism with a low inhalational infectious
dose, a concentration below the detection limit of a culture-based air sampler is not
necessarily a concentration below the infectious dose. ``Nothing was recovered'' and
``nothing was released'' are different statements, and only the first is supported.

The same caution applies in the other direction to the quantitative comparison for
centrifugation. Measured source intensities place centrifugal tube rupture at roughly
10--100~CFU/m\textsuperscript{3}, \emph{below} the roughly
10\textsuperscript{3}--10\textsuperscript{4}~CFU/m\textsuperscript{3} measured for dropped
containers and accidental injection~\cite{liu2023}. An earlier draft of this work implied
the reverse ranking. The failure-conditional flag on centrifugation is retained, but it
now rests on the breakage literature and on guidance rather than on a concentration
ranking the data do not support. That study also used 30~mL of concentrated suspension in
a research biosafety laboratory rather than clinical specimens, which is why the
corresponding specimen-reception finding is graded E2 with the inference stated rather
than E1.

\subsection{The MALDI-TOF failure mode is not the one usually taught}

The most actionable finding in this work concerns the trigger for enhanced precautions.

Conventional teaching holds that a small Gram-negative coccobacillus growing on blood and
chocolate agar should raise suspicion of a biological threat agent. An earlier version of
the control set for MALDI-TOF target preparation accordingly recommended a Gram-stain
prompt firing on that morphology.

The incident data do not support this control. In three of the four exposure events
involving the identification platform, the initial Gram stain appearance was of
Gram-positive rods or cocci suggestive of streptococci, not of small Gram-negative
coccobacilli~\cite{ackelsberg2020}. The staff looked at the Gram stain, and the Gram stain
misled them. A morphology-triggered prompt would have failed to fire in three of four
events, and would have delivered a false reassurance precisely where reassurance was most
costly.

The control set was accordingly rebuilt around a trigger that does not depend on
morphology: any slow-growing blood culture isolate is handled in a cabinet and entered
into the rule-out algorithm irrespective of its Gram stain appearance, supplemented by a
clinical and travel history cue available at the bench, and by a prompt tied to specimen
type and growth rate rather than to appearance. This aligns with sentinel-laboratory
guidance that commercial identification systems should not be used for potential select
agents, and that manipulation of organisms that might be \emph{Brucella} should occur in a
cabinet under enhanced practices until ruled out and referred~\cite{asmbruc2025}.

That the exposures also concentrate in this step --- four events accounting for 84\% of
personnel exposures identified across ten
laboratories~\cite{ackelsberg2020} --- makes the trigger design, rather than the
containment hardware, the highest-leverage change available in this part of the workflow.

\subsection{Research priorities}

Four measurement studies would resolve most of the uncertainty documented here, in
descending order of value:

\begin{enumerate}
\item \textbf{A high-volume automated analyser}, characterising normal running separately
      from cover-open fault clearance and maintenance. This is the highest-throughput step
      in the laboratory and is graded E5. A single well-designed study would move it from
      E5 to E1 and would be directly citable by every non-waived laboratory.
\item \textbf{Open-bench subculture and plate manipulation}, the step with the strongest
      incident evidence and no measurement. Measurement here would either corroborate an
      inference currently resting entirely on epidemiology, or substantially revise it.
\item \textbf{Failure modes rather than procedures}, for centrifugation and specimen
      reception. For both steps the risk is concentrated in breakage and leakage, and
      measuring the intact procedure --- as has been done --- answers the less important
      question.
\item \textbf{Decontamination adequacy of instrument internal pathways} after a
      late-recognised isolate has passed through an automated platform. Retrospective
      recognition is the norm in clinical laboratories, and the decision to return an
      instrument to service is currently made with no evidence base.
\end{enumerate}

\subsection{Relation to ISO 35001:2019}

The instrument is designed to produce the documented outputs that clauses 6.1.1 to 6.1.4
require, at a resolution the standard does not itself specify. ISO 35001 obliges risk
assessment but does not prescribe a methodology, matrix dimensionality, or whether
assessment should be agent-based or activity-based. This reference offers one
activity-based instantiation, with the clause mapping released so that a laboratory can
show which requirement each recorded score serves.

One consequence deserves emphasis. For a step graded E5, a laboratory following this
reference records explicitly that it is operating on precaution rather than evidence.
That record is itself a legitimate output under clause 6.1.4, and is more defensible than
a confident score with no provenance.

\subsection{Limitations}
\label{sec:limitations}

\textbf{These scores are structured expert judgement, not measurement.} The instrument
makes judgement explicit, decomposed and arguable; it does not make it empirical. A reader
who disagrees can identify the specific anchor applied and contest it, which is the
intent, but the figure remains a judgement.

\textbf{Inter-rater agreement has not been measured.} Every score reported here reflects
one rater's reading of the anchors. A protocol has been registered in advance, specifying
Gwet's AC2 with ordinal weights as the primary statistic and linearly weighted Cohen's
kappa as secondary, with analysis code written and self-tested against boundary cases. It
has not been executed because no second rater has yet been recruited, and reporting an
agreement coefficient without one would be fabrication. Until it is run, the instrument
should be regarded as proposed rather than validated.

\textbf{Mixed-direction anchors are a defect generator.} Dimension A contains one
reverse-coded item among five. In an earlier draft this item was aggregated without
transformation, so that better containment raised the release score; five of seven steps
were affected, and the error survived because each anchor reads correctly in isolation.
The remedy adopted --- an explicit direction note on every reverse-anchored item, and a
machine-readable declaration of reverse-coded items consumed by the aggregation --- is
recommended to anyone building a similar instrument. The defect was detected by manual
verification, not by the automated checks, which is itself a finding about the limits of
automated quality control.

\textbf{The composition parameters are asserted.} The control credit $\kappa = 0.8$ and the
band cut points at 4, 9, 14 and 19 are author-set. Sensitivity of the step ordering to
these values has not been tested.

\textbf{The literature search was targeted, not systematic.} No registered search strategy,
database list, or screening log was maintained. This release therefore does not claim
completeness of the evidence base, and a systematic search is the first planned addition.

\textbf{Baseline control scores characterise, they do not estimate.} They describe common
practice as reported in the literature and guidance. No survey of laboratories was
conducted, and the baseline for any particular laboratory may differ substantially.

\textbf{Three claims were excluded for want of a primary source}, as described in the
Methods. Their exclusion makes one control set more cautious than the full evidence may
warrant, and leaves the cabinet certification standard and interval to be recorded
locally rather than specified here.

\textbf{Agent profiles are stylised}, and scope is limited to US CLIA moderate- and
high-complexity practice. The instrument should transfer to other regulatory settings; the
baseline control scores almost certainly do not.

\section{Conclusion}

Bench biosafety practice in clinical diagnostic laboratories is inverted with respect to
its evidence base. The steps most often named in requests for containment have been
measured, with containment in place, and yielded nothing recoverable. The steps associated
with documented occupational infection and death have never been measured at all. The
busiest step in the laboratory has no evidence of any kind behind its safety practice.

This inversion is invisible to any instrument that returns a single number. Reporting the
score and the strength of the evidence beside one another, and refusing to combine them,
is what makes it legible --- and makes it possible for a laboratory to distinguish the
controls it maintains because they work from those it maintains because they always have.

The instrument, the reference scores, the extracted evidence with provenance, the control
sets, and the analysis code are released openly so that each component can be contested on
its own terms.

\section*{Data and code availability}

All data, the scoring instrument, extracted evidence with source provenance, control sets,
analysis code, and quality-control reports are openly available at
\url{https://github.com/DrAdedamola/clinical-lab-aerosol-biorisk-reference} and archived
with a persistent identifier. The instrument, data and documents are released under
CC BY 4.0; the code under the MIT licence. ISO 35001:2019 clause numbers and titles are
reproduced as a structural index only; the standard itself must be obtained from ISO or a
national member body.

\section*{Competing interests}

The author declares no competing interests.

\section*{Funding}

This work received no funding.

\section*{Use of artificial intelligence}

AI assistance was used for code implementation, data wrangling, literature triage,
document formatting, and drafting mechanics. All analytic decisions, every score, and
every citation are the author's, and each citation was checked against its source by the
author. Three claims that could not be verified at source were excluded, as described in
the Methods.

\section*{Acknowledgements}

The biorisk assessment, characterization and evaluation approach applied here was learned
through the Sandia National Laboratories biorisk management workshops delivered for the
U.S. Defense Threat Reduction Agency. No workshop material is reproduced in this work.
